\chapter{Standard}

\section{Structure}

\subsection{Project Structure}

\begin{itemize}[parsep=1pt]
	\item All files relevant for synthesis shall be contained in or in a subfolder of a general \textbf{src} folder, from here on called 'source file'.
	\item All testbenches shall be contained in or in a subfolder of a general \textbf{tb} folder.
	\item All TCL scripts shall be contained in or in a subfolder of a general \textbf{scripts} folder.
	\item All steps of VHDL development shall be conducted using Vivado TCL scripts, where possible.
\end{itemize}

\subsection{File Structure}
\label{sec:FileStructure}

\begin{itemize}[parsep=1pt]
	\item All source files shall contain exactly one entity and its one architecture, except for files describing packages.
	\item All source files shall have the same name as the contained entity, followed by a \textbf{.vhd} file extension.
	\item All source files shall have a header following Example \ref{fig:header}.
\end{itemize}
\begin{figure}[hbt]
	\lstinputlisting[language=VHDL]{code/header.vhd}
	\caption{VHDL File Header Template}
	\label{fig:header}
\end{figure}

\section{Naming Convention}
\label{sec:Naming_Convention}

\begin{itemize}[parsep=1pt]
	\item All identifiers shall be named in lower snake case, if not specified otherwise.
	\item All architectures performing clocked logic shall be named behavioral.
	\item All architectures integrating instances of entities shall be named connectivity.
	\item All signals shall be named in lower snake case.
	\item Every signal not contributing to a bus interface (e.g. AXI-Stream) shall have one of the following suffixes, depending on the signal type: 
	\begin{itemize}
		\item \textbf{\_i} for entity inputs
		\item \textbf{\_o} for entity outputs
		\item \textbf{\_r} for architecture internal signals, that get written to in a clocked process (from here on called 'registers')
		\item \textbf{\_s} for all other architecture internal signals
	\end{itemize}
	\item All constants shall be named in caps lock with underscores to delimit words (EXAMPLE\_CONSTANT).
	\item All generics shall follow the constants' naming scheme.
	\item All custom types shall have a \textbf{\_t} suffix.
	\item All identifiers should be self explanatory.
\end{itemize}
\begin{figure}[hbt]
	\lstinputlisting[language=VHDL]{code/namingConventionExample.vhd}
	\label{fig:naming}
	\caption{VHDL Example Entity}
\end{figure}

\section{Code Formatting}

\begin{itemize}[parsep=1pt]
	\item Tabs shall be used as indentations.
	\item All VHDL keywords shall be used in lower case.
	\item There shall only be one signal definition per line.
	\item All entity port lists shall be indented as in Example \ref{fig:naming}.
\end{itemize}

\section{Design Guidelines}
\label{sec:DesignGuidelines}

\begin{itemize}[parsep=1pt]
	\item All input and output signals shall be of type \textbf{std\_logic}, \textbf{std\_logic\_vector} or of a custom type.
	\item For entitys exclusive to a module, which are not exposed to the top-level or a layer's top-level, the above restriction does not apply.
	\item All maths operations shall be performed, using \textbf{signed} or \textbf{unsigned} types, where possible.
	\item Math operations should not be performed with \textbf{integers}, except for array indices.
	\item An architecture that instantiates components shall not describe clocked logic.
	\item An architecture that instantiates components should only perform minimal sequential logic, like multiplexing.
	\item All registers shall be initialized with a predefined value.
	\item There shall be no unnamed constants (so called magic-numbers) in functional code, except for logic level checks (i.e. \textbf{signal\_s = '0'}).
	\item All entities, architectures and processes shall be explicitly named.
	\item All named blocks shall be ended with its name explicitly mentioned.
	\item All code blocks shall be indented once more than their parent blocks.
	\item Every entity with a behavioral architecture shall define an active-low reset pin, which resets the entity asynchronously.
\end{itemize}

\begin{figure}[hbt]
	\lstinputlisting[language=VHDL]{code/exampleProcess.vhd}
	\label{fig:process}
	\caption{VHDL Example Process}
\end{figure}

\newpage